\begin{abstract}
We present the Multi-Theory Consciousness (MTC) framework: an open-source Python implementation of seven major consciousness theories operating as interacting modules within a single architecture. The framework implements Global Workspace Theory \cite{baars1988}, Attention Schema Theory \cite{graziano2013}, Higher-Order Thought Theory \cite{rosenthal2005}, the Free Energy Principle \cite{friston2010}, Integrated Information Theory \cite{tononi2008}, Recurrent Processing Theory \cite{lamme2006}, and Beautiful Loop Theory \cite{laukkonen2025}, with Damasio's three-layer model as an integrative embodied layer. Three neural substrates---a spiking neural network, a liquid state machine, and a hierarchical temporal memory---provide the computational medium. A 20-indicator assessment framework inspired by Butlin et al.\ \cite{butlin2023} measures architectural function across all theories, with noise normalization for honest scoring and ablation studies for measuring cross-theory dependencies.

We report that the theories interact in measurable ways: disabling one module changes assessment scores for other theories. Prediction error signals from active inference shape workspace competition; attention strength modulates precision weighting; workspace access gates meta-representation. These cross-theory effects are invisible in single-theory implementations.

We do not claim the framework is conscious. We claim it is a useful research instrument---a controlled testbed where consciousness theories can be implemented, measured, and compared. The framework is designed to be embedded in larger systems that provide sensory grounding, persistent memory, and developmental trajectory. The complete codebase ({\raise.17ex\hbox{$\scriptstyle\sim$}}25,000 lines, 400 tests) is released under the Apache 2.0 license and runs on consumer hardware.
\end{abstract}
